\documentclass[UTF8,zihao=5,fontset=windows,a4paper]{ctexart}
\usepackage[left=20mm,right=20mm,top=24mm,bottom=22mm,headheight=20pt,headsep=5mm,footskip=10mm]{geometry}
\usepackage{fontspec,xcolor,graphicx,tikz,booktabs,tabularx,array,colortbl,caption,fancyhdr,lastpage,enumitem,needspace}
\usepackage{amsmath}
\usepackage[hidelinks,unicode]{hyperref}
\usetikzlibrary{arrows.meta,positioning,calc,fit,backgrounds}
\setmainfont{Times New Roman}
\setsansfont{Calibri}
\setmonofont{Consolas}
\setCJKmainfont{SimSun}[BoldFont=SimHei]
\setCJKsansfont{Microsoft YaHei}
\definecolor{navy}{HTML}{244B6A}
\definecolor{blue}{HTML}{397BA7}
\definecolor{ink}{HTML}{202A33}
\definecolor{muted}{HTML}{536473}
\definecolor{rulegray}{HTML}{BDCEDA}
\definecolor{pale}{HTML}{EDF4F9}
\definecolor{teal}{HTML}{3D7066}
\definecolor{paleteal}{HTML}{ECF4F1}
\color{ink}
\hypersetup{pdftitle={Web Agent 的并行执行与容错调度},pdfauthor={陈思旭，寇之洲，谢树韬},pdfsubject={勇攀高峰：方案与工程过程}}
\ctexset{
 section={format=\heiti\color{navy}\fontsize{12}{18}\selectfont,number=\arabic{section},aftername=\quad,beforeskip=14pt,afterskip=8pt},
 subsection={format=\heiti\color{navy}\fontsize{10.5}{16}\selectfont,number=\thesection.\arabic{subsection},aftername=\quad,beforeskip=10pt,afterskip=6pt}
}
\setlength{\parindent}{2em}
\setlength{\parskip}{3pt}
\linespread{1.22}
\clubpenalty=10000
\widowpenalty=10000
\pagestyle{fancy}
\fancyhf{}
\fancyhead[L]{\footnotesize\color{muted} Web Agent 的并行执行与容错调度}
\fancyhead[R]{\footnotesize\color{muted} 勇攀高峰}
\fancyfoot[C]{\footnotesize\color{muted}\thepage}
\renewcommand{\headrulewidth}{0.35pt}
\renewcommand{\footrulewidth}{0pt}
\renewcommand{\headrule}{\hbox to\headwidth{\color{rulegray}\leaders\hrule height \headrulewidth\hfill}}
\captionsetup{font=small,labelfont={bf,color=navy},textfont={color=muted},labelsep=quad,skip=7pt,hypcap=false}
\DeclareCaptionType{algorithm}[算法]
\setlength{\textfloatsep}{14pt}
\setlength{\intextsep}{13pt}
\setlength{\floatsep}{13pt}
\renewcommand{\topfraction}{0.85}
\renewcommand{\bottomfraction}{0.75}
\renewcommand{\textfraction}{0.12}
\setlength{\tabcolsep}{5pt}
\newcolumntype{Y}{>{\raggedright\arraybackslash}X}
\newcolumntype{L}[1]{>{\raggedright\arraybackslash}p{#1}}
\newcommand{\kw}[1]{\textsf{\color{navy}#1}}
\newcommand{\citepaper}[1]{\textsuperscript{\textcolor{blue}{\cite{#1}}}}
\newcommand{\algline}[2]{\makebox[1.4em][r]{\scriptsize\color{muted}#1}\hspace{0.7em}#2\par}
\tikzset{
 box/.style={draw=rulegray,fill=white,line width=.6pt,align=center,inner sep=6pt,font=\small,minimum height=14mm},
 core/.style={box,draw=blue,fill=pale},
 arr/.style={-{Stealth[length=2mm,width=1.4mm]},draw=blue,line width=.75pt},
 note/.style={font=\scriptsize,text=muted,align=center}
}
\begin{document}
\thispagestyle{plain}
\begin{center}
{\heiti\color{navy}\fontsize{18}{25}\selectfont Web Agent 的并行执行与容错调度}\\[10pt]
{\fontsize{11}{16}\selectfont 陈思旭\quad 寇之洲\quad 谢树韬}\\[3pt]
{\small 第 7 队\quad 勇攀高峰}\\[4pt]
{\footnotesize 代码仓库：\href{https://github.com/xie-st/WebAgentNew}{\texttt{github.com/xie-st/WebAgentNew}}}
\end{center}
\vspace{3pt}
\begingroup
\small
\noindent\textbf{摘要：}针对网页任务中的连续交互、结果确认和容错调度问题，本文基于 EffGen 构建统一的网页智能体执行框架，通过可暂停的生成器连接官方评测平台与 Playwright 浏览器。最终官方配置采用单 Agent 串行执行，启用任务状态工具与基于页面证据的完成核验，将不合法动作和核验缺口反馈给同一执行循环。恢复与接管、必要上下文保护、批准文本经验和只读并行调度作为独立配置或扩展模块保留。工程过程由人工规划、文档约束、智能体分工与统一集成组成，形成可追溯的代码、配置和运行记录。
\par\vspace{3pt}
\noindent\textbf{关键词：}网页智能体；结果核验；状态管理；容错调度；并发控制
\par
\endgroup
\vspace{4pt}
\section{问题定义与基础工具}

本项目面向网页任务的并行执行、结果确认、容错恢复与频率控制。比赛运行模型为 Qwen3.6-35B-A3B，官方端通过 SDK 获取页面观察、调用统一模型并返回单个网页动作，真实端由 Playwright 驱动浏览器渲染、操作与截图。系统以 EffGen 为基础，复用其观察—思考—工具调用—反馈循环，以及轨迹、输入预算和历史压缩机制。\citepaper{effgen}

开发工具采用 Codex、WorkBuddy 和 DeepSeek Harness（DSH），提供文件读写、代码检索、命令执行、测试和任务派发能力。Codex 使用 GPT-6.1 Sol，推理强度为 Medium；WorkBuddy 使用 GLM-5.3-flash；DSH 经硅基流动 API 调用 DeepSeek v4 flash 1013。这些开发工具提供通用工程能力，不包括本项目所需的 Web Agent 具体知识；赛题理解、论文选择与路线规划由团队完成。

\section{工程过程与分工}

团队先确定赛题约束与优化方向，查询状态维护、观察压缩、完成核验和并发调度等论文，形成 RoadMap。Codex 将其翻译为阶段方案、接口约定、文件归属和验收要求；人工审定后，WorkBuddy 向 DSH 派发任务，主线程统一集成并核对实际入口。工程先接通两套环境，再建设状态、核验、恢复与调度能力，最后按主树入口和配置收敛提交方案。

\begin{figure}[htbp]
\centering
\begin{tikzpicture}[x=1cm,y=1cm]
\node[core,text width=3.0cm] (a) at (0,0) {\textbf{团队}\\[-1pt]\footnotesize 赛题 · 论文 · RoadMap};
\node[box,text width=3.0cm] (b) at (4.2,0) {\textbf{Codex}\\[-1pt]\footnotesize 方案 · 实施 · 验收};
\node[box,text width=3.0cm] (c) at (8.4,0) {\textbf{WorkBuddy}\\[-1pt]\footnotesize 编排与派发};
\node[core,text width=3.0cm] (d) at (12.6,0) {\textbf{DSH}\\[-1pt]\footnotesize 实现与离线验证};
\draw[arr] (a)--(b);
\draw[arr] (b)--(c);
\draw[arr] (c)--(d);
\draw[arr] (d.south)--++(0,-.45)-| (a.south);
\node[note,fill=white,inner sep=2pt] at (6.3,-.93) {主线程集成与实际入口验证\quad → \quad 团队反馈与下一批决策};
\end{tikzpicture}
\caption{人机协作的工程开发流程}\label{fig:workflow}
\end{figure}

并行开发沿用模块边界。输入与观察适配、状态与恢复支持、本地任务与判定分别在独立源码副本中实现；运行测试前确认导入各自副本。公共驱动、官方入口、配置和命令行由一名负责人顺序合入，跨模块功能在接口稳定后连接。真实模型与官方平台任务顺序执行，开发智能体的并行不改变比赛执行端的并发约束。

\Needspace{7\baselineskip}
\section{最终官方执行方案}

\subsection{统一接口与可暂停循环}

最终入口为 \path{webagent.official_agent.OfficialAgent}，提交配置为 \path{code/harness/configs/official.yaml}。每个任务建立独立 Session、状态库、模型适配器和 EffGen 驱动。官方端模型调用仅经过 SDK；统一网页工具向后端交付点击、输入、按键、滚动、导航等动作，模型依据返回的可访问性树（AX）继续决策。

\begin{figure}[htbp]
\centering
\begin{tikzpicture}[x=1cm,y=1cm]
\node[core,text width=4.2cm,minimum height=18mm] (agent) at (0,0) {\textbf{EffGen 决策循环}\\[-1pt]\footnotesize Qwen3.6-35B-A3B};
\node[box,text width=3.0cm,minimum height=18mm] (tools) at (5.2,0) {\textbf{统一网页工具}\\[-1pt]\footnotesize 动作、观察与元素语义};
\node[box,text width=3.2cm] (official) at (10.2,.9) {\textbf{官方平台}\\[-1pt]\footnotesize SDK · 单动作交互};
\node[box,draw=teal,fill=paleteal,text width=3.2cm] (real) at (10.2,-.9) {\textbf{真实浏览器}\\[-1pt]\footnotesize Playwright · 渲染与截图};
\draw[arr] ([yshift=3mm]agent.east)--node[above,note]{交出动作}([yshift=3mm]tools.west);
\draw[arr] ([yshift=-3mm]tools.west)--node[below,note]{接回观察}([yshift=-3mm]agent.east);
\draw[arr] (tools.east)--++(.35,0)|- (official.west);
\draw[arr] (tools.east)--++(.35,0)|- (real.west);
\node[note,draw=rulegray,fill=pale,rounded corners=1pt,text width=10.1cm,minimum height=8mm] (evidence) at (3.8,-2.1) {状态与证据：目标 · 观察 · 动作反馈 · 完成核验 · 模型调用 · 截图};
\draw[draw=rulegray,line width=.6pt] (agent.south)--(agent.south|-evidence.north);
\draw[draw=rulegray,line width=.6pt] (tools.south)--(tools.south|-evidence.north);
\end{tikzpicture}
\caption{共用决策循环的双环境接入结构}\label{fig:architecture}
\end{figure}

网页动作处通过生成器暂停，下一次 \path{step()} 将观察送回原循环；状态读取、计划更新和证据记录等内部工具直接本地执行。动作绑定页面、观察版本和当前有效元素编号，错误目标或不合法参数在发送前拒绝，并通过工具反馈继续纠正，默认最多允许两次格式纠正。

输入采用 \path{decompose} 语义：多行文本拆成输入与 Enter 子动作，每个子动作接收新的页面反馈，完整序列结束后再交回模型；空文本提交转换为 Enter，空文本且不提交的动作直接反馈纠正。文本替换由明确的选中、清除和输入动作完成，保留真实浏览器输入语义。


\begin{algorithm}[htbp]
{\color{blue}\hrule height .6pt}
\caption{最终官方入口的连续执行流程}\label{alg:loop}
{\color{rulegray}\hrule height .3pt}
\vspace{6pt}
\begingroup
\small\setlength{\parindent}{0pt}\setlength{\parskip}{3pt}
\algline{1}{建立任务预算、状态库和会话，注册启用工具，安装首个观察}
\algline{2}{\kw{while} 任务活动且未到平台截止时间：}
\algline{3}{\hspace*{1.3em}推进同一 EffGen 生成器，取得提案 $p$}
\algline{4}{\hspace*{1.3em}\kw{if} $p$ 是内部工具：本地执行并反馈；\kw{continue}}
\algline{5}{\hspace*{1.3em}\kw{if} $p$ 请求完成：核验已声明目标；通过则结束}
\algline{6}{\hspace*{2.6em}否则返回缺口继续修正，超出限额则记为 \texttt{incomplete}}
\algline{7}{\hspace*{1.3em}\kw{else} 校验观察版本与动作目标；不合法则反馈纠正}
\algline{8}{\hspace*{2.6em}必要时拆分输入；逐个交出动作，接回新观察和反馈}
\algline{9}{\hspace*{2.6em}保存动作与证据，推进原循环}
\algline{10}{保存终态、核验结果、请求记录与终止原因}
\endgroup
\vspace{2pt}
\noindent{\color{rulegray}\rule{\linewidth}{.4pt}}
\end{algorithm}


\subsection{状态账本与完成核验}

默认注册 \path{goal_declare}、状态读写和计划更新工具。模型声明业务要求及可观察的完成条件，程序绑定当前版本、分支和证据来源。状态库保存目标、事实、依赖与证据引用，更新时校验版本，再将完整 JSON 快照原子替换；截图与原始观察沿用证据存储，不在账本中重复保存。

完成判定检查已声明目标。官方路径使用确定性条件核对当前 URL、AX 文本、字段值和已保存证据，核验器自身不执行网页动作。\path{request_finish} 未通过时，把未满足的目标与条件反馈给同一生成器，允许两次继续修正，超过限额则以 \path{incomplete} 结束。字段值核验适配官方 AX 的属性后缀和子级文本，区分操作意图与实际页面结果。

\subsection{提交配置与预算}

下表依据当前主树配置及其实际解析结果。状态工具与完成核验进入默认执行路径；其余扩展通过独立开关控制。

\begin{table}[htbp]
\caption{主树最终官方配置}\label{tab:config}
\centering\begingroup
\linespread{1.02}\fontsize{9.5}{14}\selectfont
\renewcommand{\arraystretch}{1.18}
\begin{tabularx}{\linewidth}{@{}L{31mm}Y@{}}
\toprule\rowcolor{pale}
\textbf{\color{navy}项目} & \textbf{\color{navy}最终官方配置}\\
\midrule
决策与执行 & 单 Agent、文本 ReAct；同一时刻仅一个任务、一个动作和一次模型调用；关闭自动子智能体、推测执行和跨任务记忆。\\
默认能力 & \path{state_tools=true}，\path{verification=true}。\\
可选开关 & \path{recovery=false}，\path{context.enabled=false}，\path{experience.enabled=false}，\path{parallel.enabled=false}。\\
输入预算 & 原生压缩输入预算 24,000 token，规划上下文 32,768 token；采用估算计数。\\
请求约束 & 模型出口检查最多 256 条消息、1 MiB JSON 请求；必要任务与当前观察不因字节超限而静默删除。\\
截止与终态 & 提前截止余量为 0，遵守平台绝对截止时间；终止时保存结果、调用记录和错误原因。\\
\bottomrule
\end{tabularx}
\endgroup\end{table}

\Needspace{7\baselineskip}
\section{可选机制与并发调度}

\subsection{恢复、接管与上下文}

开启恢复开关后，官方入口连接故障分类、停滞检测和有界恢复。恢复依据失败证据与目标依赖撤回受影响的完成判断，保留其他已确认进度；每目标最多两次、全任务最多五次恢复。待确认副作用在动作发出前入账，反馈未知时保留未知状态，先核实再决定后续动作。

同一开关接入单任务人工接管：登录、验证或主动暂停时关闭动作出口，处理已执行反馈后等待明确继续信号；恢复时重新观察，废弃旧提案与旧元素编号。快照重启恢复业务状态与证据，并重新建立执行循环。

必要上下文保护开启后，在模型发送前核对任务、当前 AX、账本短视图及待确认副作用，超限时明确拒绝发送。经验加载开启后，仅从本地批准清单选择版本与适用条件匹配的文本，进入既有提示词装配；当前提交配置关闭这两项，默认执行沿用 EffGen 原生上下文管理。

\subsection{只读并行与统一限流}

调度模块通过独立协调器接收分支任务，每个分支关联会话、循环和页面槽位；仅独立只读分支可并行，计划中出现写任务时整批退回串行。所有分支共用模型调用门控、访问节奏与根预算，失败、取消和丢弃分支同样计入汇总。模型并发固定为一，站点操作间隔默认一秒，遇到冷却提示时共享退避状态。

当前正式官方入口采用串行路径，配置拒绝开启多 Tab；独立调度模块保留给真实端扩展。并行模块的多分支暂停、继续与排空仍是预留接口，单任务恢复与接管由官方入口中的控制层承担。

\subsection{论文方法与方案取舍}

论文为机制选择提供依据，具体接入以当前代码和配置为准。状态、核验与统一执行进入官方主流程，其他方法按下表保留相应模块或候选形式。

\begin{table}[htbp]
\caption{论文方法与最终方案的对应}\label{tab:methods}
\centering\begingroup
\linespread{1.02}\fontsize{9.5}{14}\selectfont
\renewcommand{\arraystretch}{1.18}
\begin{tabularx}{\linewidth}{@{}L{25mm}L{39mm}Y@{}}
\toprule\rowcolor{pale}
\textbf{\color{navy}方向} & \textbf{\color{navy}论文依据} & \textbf{\color{navy}本项目采用方式}\\
\midrule
状态与停滞 & PoS / BeliefState \citepaper{pos} & 账本区分事实和目标缺口；停滞诊断随恢复开关启用。\\
核验与纠错 & WALT \citepaper{walt}、Reflexion \citepaper{reflexion} & 默认核验已声明目标，反馈送回同一循环；有界恢复按证据撤回受影响状态。\\
摘要与经验 & M² \citepaper{m2} & 提供批准文本经验加载，默认关闭；同轮摘要保留为候选。\\
历史管理 & AgentProg \citepaper{agentprog}、Read More, Think More \citepaper{readmore} & 复用原生轨迹与压缩；子目标折叠、历史差分保留为候选。\\
页面筛选 & Prune4Web \citepaper{prune} & 关键词与固定程序评分保留为候选，当前主流程读取 AX。\\
依赖调度 & Flash-Searcher \citepaper{flash} & 独立协调器实现只读分支调度、共享写退回串行及根预算汇总。\\
视觉与行动效率 & UI-TARS \citepaper{uitars}、\mbox{WebClipper\citepaper{clipper}} & 截图用于留证；截图参与决策、轨迹裁剪及动作组合保留为候选。\\
\bottomrule
\end{tabularx}
\endgroup\end{table}

\section{工程交付与追溯}

提交工程包含 \path{code/harness/}、固定版本及补丁的 EffGen、官方 SDK 和配置。两套后端共用观察、动作、会话与证据协议；日志关联任务、页面、动作反馈、模型请求和终止原因。官方请求的真实 token 用量以平台记录为准，本地估算值单独标明。最终方案以本次读取的主树源码和配置为依据，相关文件哈希保存在报告来源目录，便于材料与工程版本核对。

\begingroup
\renewcommand{\refname}{参考文献}
\ctexset{section={beforeskip=5pt,afterskip=4pt}}
\linespread{1.0}
\fontsize{9.3}{13.2}\selectfont
\setlength{\parskip}{0pt}
\begin{thebibliography}{11}
\raggedright
\setlength{\itemsep}{3pt}
\setlength{\parsep}{0pt}
\bibitem{effgen} \href{https://arxiv.org/abs/2602.00887}{EffGen: Enabling Small Language Models as Capable Autonomous Agents}. arXiv:2602.00887, 2026.
\bibitem{pos} \href{https://arxiv.org/abs/2610.01415}{Beyond Memory: Harnessing Long-Horizon Agents with Explicit Belief States}. arXiv:2610.01415, 2026.
\bibitem{walt} \href{https://arxiv.org/abs/2510.01524}{WALT: Web Agents that Learn Tools}. arXiv:2510.01524, 2025.
\bibitem{reflexion} \href{https://arxiv.org/abs/2303.11366}{Reflexion: Language Agents with Verbal Reinforcement Learning}. arXiv:2303.11366, 2023.
\bibitem{m2} \href{https://arxiv.org/abs/2603.00503}{M²: Dual-Memory Augmentation for Long-Horizon Web Agents via Trajectory Summarization and Insight Retrieval}. arXiv:2603.00503, 2026.
\bibitem{agentprog} \href{https://arxiv.org/abs/2512.10371}{AgentProg: Empowering Long-Horizon GUI Agents with Program-Guided Context Management}. arXiv:2512.10371, 2025.
\bibitem{readmore} \href{https://arxiv.org/abs/2604.01535}{Read More, Think More: Revisiting Observation Reduction for Web Agents}. arXiv:2604.01535, 2026.
\bibitem{prune} \href{https://arxiv.org/abs/2511.21398}{Prune4Web: DOM Tree Pruning Programming for Web Agent}. arXiv:2511.21398, 2025.
\bibitem{flash} \href{https://arxiv.org/abs/2509.25301}{Flash-Searcher: Fast and Effective Web Agents via DAG-Based Parallel Execution}. arXiv:2509.25301, 2025.
\bibitem{uitars} \href{https://arxiv.org/abs/2501.12326}{UI-TARS: Pioneering Automated GUI Interaction with Native Agents}. arXiv:2501.12326, 2025.
\bibitem{clipper} \href{https://arxiv.org/abs/2602.12852}{WebClipper: Efficient Evolution of Web Agents with Graph-based Trajectory Pruning}. arXiv:2602.12852, 2026.
\end{thebibliography}
\endgroup
\end{document}
